\section*{第 \ref{ch:b1:ecosystem-dynamics} 章——生态系统动态}

\begin{solution}{exo:b1:ecosystem-dynamics:1}
\omterm{def:b1:ecosystem-dynamics:succession}{初生演替}：在没有土壤的新基质上发生（如冰碛、熔岩流）；
\omterm{def:b1:ecosystem-dynamics:succession}{次生演替}：土壤与种子库尚存之处（如废弃农田、
烧毁的林地）。\omterm{def:b1:ecosystem-dynamics:succession}{次生演替}快得多，约快十倍或更多，原因在于
土壤、养分与种子都无需重新形成。
\end{solution}

\begin{solution}{exo:b1:ecosystem-dynamics:2}
促进、抑制、耐受。在熔岩上属促进作用：先锋物种
（\omterm{def:b1:species-interactions:symbiosis}{地衣}、苔藓、\omterm{def:b1:plant-water-minerals:nodules}{固氮}生物）须先形成土壤，
其他物种才能扎根。
\end{solution}

\begin{solution}{exo:b1:ecosystem-dynamics:3}
\omterm{def:b1:ecosystem-dynamics:disturbance}{抗性}：在\omterm{def:b1:ecosystem-dynamics:disturbance}{干扰}中保持不变；\omterm{def:b1:ecosystem-dynamics:disturbance}{恢复力}：
即\omterm{def:b1:ecosystem-dynamics:disturbance}{干扰}后迅速恢复。珊瑚礁对变暖既无\omterm{def:b1:ecosystem-dynamics:disturbance}{抗性}，
亦无\omterm{def:b1:ecosystem-dynamics:disturbance}{恢复力}；草甸对火的\omterm{def:b1:ecosystem-dynamics:disturbance}{抗性}低而
\omterm{def:b1:ecosystem-dynamics:disturbance}{恢复力}高。
\end{solution}

\begin{solution}{exo:b1:ecosystem-dynamics:4}
$16^{0.25} = 2$：\num{24}~个物种；$625^{0.25} = 5$：\num{60}~个物种。
\end{solution}

\begin{solution}{exo:b1:ecosystem-dynamics:5}
冰碛土几乎不含氮；云杉不能\omterm{def:b1:plant-water-minerals:nodules}{固氮}，只在
有氮供应的地方才能生长。赤杨能\omterm{def:b1:plant-water-minerals:nodules}{固氮}，使
土壤含氮量升至每公顷数百千克，因此
正是它使该地点适合云杉生长：促进作用。
\end{solution}

\begin{solution}{exo:b1:ecosystem-dynamics:6}
$S^* = I_0 P/(I_0 + eP) = 1600/(8 + 10) = 89$ 个物种；周转率
$eS^* = 4.4$ 个物种/年。取 $I_0 = 4$ 时，$800/14 = 57$ 个物种，
周转率为 \num{2.9}。
\end{solution}

\begin{solution}{exo:b1:ecosystem-dynamics:7}
\omterm{def:b1:ecosystem-organization:trophic}{分解者}继续将\omterm{def:b1:ecosystem-dynamics:decomposition}{凋落物}和腐殖质矿化为铵，
硝化细菌再将其转化为硝酸根，这是一种土壤无法保持的可溶性阴离子，
因无根系吸收，它被冲走。再生植被
恢复了吸收：新植物吸收硝酸根的速度与其生成速度一样快，
溪流重新变清。
\end{solution}

\begin{solution}{exo:b1:ecosystem-dynamics:8}
$k = \ln 2/1 = \qty{0.69}{yr^{-1}}$；三年后剩余 $e^{-2.08} =
\qty{12.5}{\%}$。针叶树：$k = -\ln 0.85 = \qty{0.16}{yr^{-1}}$，
半衰期 \num{4.3} 年，三年后剩余 \qty{61}{\%}——
\omterm{def:b1:ecosystem-dynamics:decomposition}{凋落物}不断积累。
\end{solution}

\begin{solution}{exo:b1:ecosystem-dynamics:9}
未放牧：高大的\omterm{def:b1:enzymes:inhibitors}{竞争性}禾草遮蔽其余植物，物种很少。
适度放牧：优势种被啃食，光线能照到
地面，小型草本与禾草共存，多样性最高。重度
放牧：只有匍匐、不可食或生长迅速的物种
能在践踏和啃食中存活，物种又变少了。
\end{solution}

\begin{solution}{exo:b1:ecosystem-dynamics:10}
原始：$c\,1000^{0.25} = 5.6c$；单个片段 $10^{0.25}c = 1.78c$，
为原始的 \qty{32}{\%}。若十个隔离片段的名录
完全不同，则合计 $17.8c$——超过原始值，
这显然荒谬：片段间共有大部分物种，因此相加
会重复计数。实际总数介于 $1.78c$（名录相同）与
约 $5.6c$（原始）之间，再减去隔离与边缘效应
所造成的损失；最先消失的，是需要超过
\qty{10}{km^2} 面积的物种。
\end{solution}

\begin{solution}{exo:b1:ecosystem-dynamics:11}
若数目完全恢复，则提示存在由岛屿条件决定的固定\omterm{def:b1:ecosystem-organization:ecosystem}{群落}；
数目可在名录不同的情况下恢复，而名录持续变化，
这正是随机迁入与灭绝之间平衡的预测——数目
是面积与距离的属性，
物种身份则纯属偶然。
\end{solution}

\begin{solution}{exo:b1:ecosystem-dynamics:12}
并非封闭：它从外界获得雨水、尘埃、大气氮（通过\omterm{def:b1:plant-water-minerals:nodules}{固氮}
和沉降），以及岩石风化释放的离子，并损失水分、
溶解离子和气体（经溪流与大气）；但
相对于从\omterm{def:b1:ecosystem-dynamics:decomposition}{凋落物}到土壤、根系再到叶片的内部循环，输入与输出都很小，
因此以一年收支的\omterm{def:b1:cell-unit-of-life:resolution}{分辨率}看，森林似乎是封闭的。砍伐
揭示了这种表观封闭的来源：
活体植被的吸收，植被一旦去除，
同样的土壤和\omterm{def:b1:ecosystem-organization:trophic}{分解者}使山谷
在一个季节内流失养分。
\end{solution}

\begin{solution}{pb:b1:ecosystem-dynamics:1}
\textbf{1.} 10--40 年：$140/30 = \qty{4.7}{kg/ha}$/年；40--80：
$150/40 = \num{3.75}$；80--200：$20/120 = \num{0.17}$。最快出现在
赤杨阶段。
\textbf{2.} 从赤杨到云杉的促进作用：赤杨所添加的氮
正是云杉所需要的。
\textbf{3.} 森林已达到\omterm{def:b1:organism-environment:homeostasis}{稳态}：输入（\omterm{def:b1:plant-water-minerals:nodules}{固氮}、
沉降）与损失（淋溶、反硝化、木材移出）相平衡，
氮在系统内部循环，而非不断积累。
\textbf{4.} 40 年间保留 \qty{150}{kg/ha}，相对于
固定的 \num{2000}：为 \qty{7.5}{\%}；其余保留在赤杨自身
的生物量与\omterm{def:b1:ecosystem-dynamics:decomposition}{凋落物}中，或被淋溶、反硝化。
\textbf{5.} 在暴露 20 年的冰碛土上，设有赤杨与无赤杨、
施氮与不施氮的样地，种植云杉幼苗。若云杉在无赤杨且
施肥时生长、不施肥则不生长，则赤杨的作用是供氮（促进作用）；
若云杉在各样地生长相同，
只是比赤杨慢，则为耐受。
\textbf{6.} 这一推断假定各地点起点相同、沿同一路径发展，
而气候、偶然迁入与冰川冰碛土的差异可能打破这一路径；
只有永久样地才能在同一
地点证明这一\omterm{def:b1:ecosystem-dynamics:succession}{演替}序列。
\textbf{7.} 到 \qty{330}{mm}，\qty{40}{\%}：树木
已不再蒸腾其应占的降雨份额。
\textbf{8.} 对照区：\qty{2}{kg/ha} 在 \qty{8200}{m^3} 中：\qty{0.24}{mg/L}。
砍伐区：\qty{120}{kg} 在 \qty{11500}{m^3} 中：\qty{10.4}{mg/L}，
为对照的 40 倍。
\textbf{9.} 对照区：$6.5 - 2 = +\qty{4.5}{kg/ha}$ 每年，逐年积累。
砍伐区：$6.5 - 120 = -\qty{113.5}{kg/ha}$ 每年，持续流失。
\textbf{10.} $120/3850 = \qty{3.1}{\%}$ 每年；该库
按此速率约可维持 \num{30} 年——但速率随
再生长恢复而下降。
\textbf{11.} 钙增加 10 倍，钾增加 18 倍。钾主要
保存在活\omterm{def:b1:body-plans-tissues:tissue}{组织}中，且易从\omterm{def:b1:ecosystem-dynamics:decomposition}{凋落物}中淋失；钙
还结合在土壤黏土与岩石中，释放较慢。
\textbf{12.} 来源于\omterm{def:b1:ecosystem-dynamics:decomposition}{凋落物}、腐殖质及砍伐残体的矿化：
先由\omterm{def:b1:ecosystem-organization:trophic}{分解者}进行氨化作用（有机氮转化为铵），
再由硝化细菌进行硝化作用（铵转化为硝酸根），在
更温暖、湿润且无覆盖的土壤中更快。
\textbf{13.} 由 \num{120} 降至 \num{20}，再降至 \num{5}：
约第 6 年时低于 \qty{6.5}{kg/ha} 的输入量；新植被
（宾州樱桃、覆盆子及幼树）的吸收恢复了养分保留。
\textbf{14.} $\log A$：0, 1, 2, 3, 4；$\log S$：1.08, 1.32, 1.58，
1.83, 2.08——为一条直线。
\textbf{15.} $z = (2.08 - 1.08)/4 = 0.25$；$c = 12$。中间岛屿：
$12\times 10^{0.25} = 21.3$，$12\times 10^{0.5} = 37.9$，$12\times
10^{0.75} = 67.5$：相差均不超过一个物种。
\textbf{16.} $12\times 250^{0.25} = 12\times 3.98 = 48$ 个物种。
\textbf{17.} $38 = 6\times 300/(6 + 300e)$：$6 + 300e = 47.4$，$e =
\qty{0.138}{yr^{-1}}$；周转率 $0.138\times 38 = 5.2$ 个物种/年。
\textbf{18.} $3\times 300/(3 + 41.4) = 20$ 个物种。
\textbf{19.} $30\times 300/(30 + 41.4) = 126$ 个物种。灭绝
在 \qty{100}{km^2} 岛屿上仍然很快：即使迁入不受限制，
那些无法在此面积上维持\omterm{def:b1:populations:population}{种群}的物种仍将
持续灭绝；决定上限的是面积，而非迁入。
\textbf{20.} $z$ 越大，物种数随面积
减小而下降越快：小岛相对更为贫乏。从生物学看，
小岛上的\omterm{def:b1:populations:population}{种群}更容易灭绝，或岛屿隔离程度很高，
以致很少物种能抵达（隔离会提高 $z$）。
\textbf{21.} $(1000/100\,000)^{0.25} = 0.01^{0.25} = 0.32$：约为
三分之一。
\textbf{22.} 隔离时，每个保护区保有 $(100/100\,000)^{0.25} = 0.18$ 的
物种；若各自名录彼此独立，则它们合计
可多于单一保护区，但彼此重叠很多，且各自都会失去
需要超过 \qty{100}{km^2} 面积的物种，因此实际更少。
若以廊道相连，它们对利用廊道的物种而言，相当于一个
\qty{1000}{km^2} 的保护区：物种数约为三分之一，且
具有的优势是局部灾难不会使所有物种一并消失。
\textbf{23.} \qty{1000}{km^2} 边长为 \qty{31.6}{km}；内部
面积 $(31.6 - 0.6)^2 = \qty{961}{km^2}$，\qty{96}{\%}。\qty{100}{km^2} 的
边长为 \qty{10}{km}：$(10 - 0.6)^2 = 88$，\qty{88}{\%}。
\textbf{24.} 物种数从原有数目缓慢下降，趋向新的
平衡 $cA^z$，历时数十年：注定灭绝的物种
起初仍然存在，只是\omterm{def:b1:populations:population}{种群}太小，难以维持——
灭绝债务即当前物种数与
平衡值之差，尚待偿还。
\textbf{25.} $z = 0.25$；\qty{100}{km^2} 岛屿平衡时保有 \num{38}~
个物种，连通后约 \num{126} 个。
\end{solution}
